\section{引言：为什么需要自监督学习}

上一讲我们讨论了 GPU、分布式训练和大规模模型训练的工程问题。本讲要回答的是：\textbf{如果我们有海量原始图像，但没有昂贵的人工标注，能不能先把表示学好？}

\begin{figure}[H]
\centering
\includegraphics[width=0.95\textwidth]{slides-images/slide-001.jpg}
\caption{课程标题页：Self-Supervised Learning\protect\footnotemark}
\end{figure}
\footnotetext{Slides 第1页。}

自监督学习的核心目标不是直接解决下游任务，而是先通过一个“预文本任务”（pretext task）让模型学出有用的表示，再把这些表示迁移到分类、检测、分割等下游任务。

\begin{knowledgebox}{为什么需要表示学习}
\begin{itemize}
    \item 直接监督学习需要大量人工标注，尤其是分割、检测这类任务，成本极高。
    \item 如果模型能从原始数据中学到“视觉常识”，这些表示往往可以迁移到很多任务。
    \item 评价一个自监督方法时，真正重要的通常不是预文本任务本身，而是它学到的特征对下游任务有多有用。
\end{itemize}
\end{knowledgebox}

\subsection{从特征到标签}

课程前半段我们已经看到，深层网络中间层的特征往往比原始像素更有语义。如果能学到一个好的 encoder，就可以只在少量标注数据上训练一个很浅的分类器，甚至直接做线性 probing。

\begin{figure}[H]
\centering
\includegraphics[width=0.95\textwidth]{slides-images/slide-008.jpg}
\caption{有用的中间表示可以通过线性分类器直接转成标签\protect\footnotemark}
\end{figure}
\footnotetext{Slides 第8页。}

这节课的主要思想可以概括为一句话：
\begin{importantbox}{自监督学习的基本范式}
先用无标注数据定义一个可自动生成标签的任务，训练 encoder 学表示，再把 encoder 迁移到真正关心的下游任务上。
\end{importantbox}

%% ========== 预文本任务 ==========

